%------------------------------------------------------------------------------%
\chapter{Introdução ao Pensamento Matemático}

Resolver problemas matemáticos vai além de ``fazer contas''. É necessário saber
argumentar e justificar cada passagem e o que estamos assumindo como verdade.
Sendo assim, resolver um exercício de matemática consiste na construção de um
texto argumentativo onde é fundamental que o aluno pense na resolução como uma
redação argumentativa. Pensando nisso, é inevitável que os estudantes tenham em
mente alguns conceitos referentes a lógica matemática que irão auxiliar-los na
construção do pensamento matemático.

%------------------------------------------------------------------------------%
\section{Proposição, Teoremas}

Uma \emph{Proposição} é uma oração afirmativa, cujo valor lógico é falso ou
verdadeiro, sem dar espaço para uma terceira alternativa. Considere, como
exemplo, as seguintes proposições:
\begin{enumerate}[label=\alph*)]
    \item Todo triângulo é retângulo;
    \item A soma de todos os ângulos de um triângulo é \ang{180};
    \item A soma de dois números pares é sempre um número par;
    \item Todo Brasileiro é carioca;
    \item Se $a>1$ então $a^2<a$.
\end{enumerate}
Note que as afirmações (a), (d) e (e) são falsas. Para verificar que, de
fato, que uma proposição é falsa, utilizamos o que chamamos de
\emph{contraexemplos}. \emph{Contraexemplos} são casos que contrariam uma
proposição, ou seja, um exemplo que refuta uma determinada afirmação universal.

Um triângulo cujos ângulos interno são \ang{120}, \ang{40} e
\ang{20} não é um triângulo retângulo, uma vez que não possui um ângulo
interno de \ang{90}. Logo esse contra-exemplo refuta a proposição (a)
mostrando que tal afirmação é falsa. Para verificar a veracidade da afirmação
(d) teríamos que consultar o registro de todos os brasileiros e verificar se
nasceu no Rio de Janeiro, mas isto é falso, pois o conhecido escritor Graciliano
Ramos é um brasileiro nascido em Alagoas. Do mesmo modo, se escolhermos $a=2$,
podemos ver que $a^2=4 > a$ e, portanto, esse contraexemplo contradiz a
afirmação feita em (e), mostrando que essa também é falsa.

Por outro lado, as proposições (b) e (c) são verdadeiras e para nos
convencermos precisaríamos demonstrar usando argumentos matemáticos. Adiante
falaremos mais um pouco sobre isso. Observe que a proposição (e) é do tipo
``Se $P$, então $Q$'', onde $P$ e $Q$ são sentenças. No item (e), por exemplo,
temos que
\begin{gather*}
  P\colon a>1 \\
  Q\colon a^2<a.
\end{gather*}
Isso significa que estamos assumindo que $P$ é verdade e usando este fato
devemos verificar se $Q$ é verdade ou não. Uma \emph{proposição condicional}
ou \emph{implicativa} é uma nova proposição formada a partir de duas
proposições P e Q, que é escrita na forma:
\[
  \text{Se }       P,
  \text{ então }   Q
  \text{ ou }      P
  \text{ implica } Q
\]
Denotamos a expressão
$P$ implica $Q$ por $P \Rightarrow Q$

No caso, a sentença P é chamada de \emph{hipótese} e a sentença $Q$ é
denominada de \emph{tese} e a validade da hipótese implica na veracidade da
tese.

A \emph{recíproca} de uma proposição condicional ou implicativa é a
proposição:
\[
  \text{Se }     Q,
  \text{ então } P
  \text{ ou }    Q \Rightarrow P
\]
A \emph{Contrapositiva} ou \emph{Contrarecíproca}
de uma afirmação condicional é a expressão :
\[
  \text{Se }     \sim Q,
  \text{ então } \sim P
  \text{ ou }    \sim Q
  \Rightarrow    \sim P
\]
Vale ressaltar que a contrapositiva tem o mesmo valor lógico da proposição a
ela associada, ou seja, dizer que $P \Rightarrow Q$ é o mesmo que dizer
$\sim Q \Rightarrow \sim P$.
Vamos chamar o modo em enunciamos uma proposição de
\emph{forma  positiva}.

Para que fique claro, considere a seguinte proposição:
\[
  \text{Se como laranja então gosto de fruta.}
\]
Na proposição enunciada acima temos que a hipótese é ``como laranja'' e a tese
é ``gosto de frutas''. Por outro lado a recíproca dessa proposição é:
\[
  \text{Se gosto de frutas então como laranja.}
\]

Por fim, a contrapositiva é:
\[
  \text{Se não gosto de frutas então não como laranja.}
\]

Mais a frente veremos que a forma contrapositiva de uma proposição poderá ser,
eventualmente, uma forma indireta muito eficaz de verificar resultados em
Matemática.

A teoria matemática é construída pelo \emph{Modelo Axiomático}, isso
significa que partimos de alguns conceitos primitivos e afirmações básicas
(\emph{axiomas ou postulados}) que aceitamos como verdadeiras, isto é, sem
que sejam demonstradas. Os conceitos de ponto, reta e plano são conceitos
primitivos. Quando afirmamos que ``Por dois pontos passam uma única reta'',
temos um exemplo de um axioma. A partir desses conceitos, todos os outros
conceitos podem ser definidos e todas as outras proposições demonstradas.
Assim, \emph{definições}  conceitos dados em função de termos considerados
previamente conhecidos.

Existem apenas duas categorias de enunciados: as proposições primitivas
(axiomas), que são proposições aceitas sem demonstração (não havendo
preocupação se são evidentes ou não); e as proposições demonstradas (teoremas,
proposições, corolários e lemas) por meio de raciocínios logicamente corretos,
a partir dos postulados e definições pertinentes.

Nesse sentido, entendemos por \emph{demonstração} Matemática, o processo de
raciocínio lógico e dedutivo para checar a veracidade de uma proposição
condicional. Nesse processo são usados argumentos válidos e definições, ou
seja, aqueles que concluam afirmações verdadeiras a partir de fatos que também
são verdadeiros.

As proposições que podem ser demonstradas são chamadas de \emph{Teoremas}.
Em geral, denominamos de Teorema apenas certas proposições que são de grande
importância matemática, chamando-se simplesmente de proposição ao resto das
proposições verdadeiras que admitem uma demonstração. Para que um Teorema seja
verificado, as hipóteses tem que ser cuidadosamente satisfeitas. Um famoso
Teorema é o \emph{Teorema de Pitágoras}, que enunciaremos a seguir:

\begin{teorema}{Teorema de Pitágoras}{teo:pitagoras}
  Em um triângulo retângulo a
  soma dos quadrados dos catetos é igual ao quadrado da hipotenusa.
\end{teorema}

Note que o teorema de Pitágoras não está na forma ``Se P então Q''. No entanto,
podemos reescrevelo nesse formato.

\begin{teorema}{Teorema de Pitágoras no Gradado}{teo:pitagoras-quadrado}
  Se $T$ é um triângulo retângulo
  de catetos $b$ e $c$ e hipotenusa $a$, então $a^2=b^2+c^2$.
\end{teorema}

No Teorema de Pitágoras temos que as hipóteses são ``T é um triângulo retângulo
de catetos $b$ e $c$ e hipotenusa $a$'' e a tese é ``$a^2=b^2+c^2$ ''. Existem
algumas sentenças que são diretas consequências de um determinado teorema são
chamadas de \emph{Corolários}, já resultados auxiliares que são usados para
demonstrar uma determinada proposição ou teorema são chamados de
\emph{Lemas}.

Quando em uma sentença temos que $P \Rightarrow Q$ e $Q \Rightarrow P$,
escrevemos $P \Leftrightarrow Q$ e o símbolo $\Leftrightarrow$ é lido como ``se
e somente se''. Em algumas proposições e teoremas matemáticos vale a recíproca,
mas nem sempre isso é verdade.

%------------------------------------------------------------------------------%
\section{Tipos de demonstração}

Como já vimos anteriormente, a demonstração  é um processo construído afim de
constatar a autenticidade de uma proposição. Durante o processo de
demonstração as hipóteses são assumidas e através de definições, deduções
feitas a partir das hipóteses ou uso de outros resultados e axiomas temos o
objetivo de chegar ao resultado, isto é, chegar a tese. Aqui serão apresentadas
3 tipos de demonstrações, a saber:
\begin{enumerate}
    \item Demonstração direta;
    \item Demonstração pela contrapositiva;
    \item Demonstração por contradição ou redução ao absurdo.
\end{enumerate}

%------------------------------------------------------------------------------%
\subsection{Demonstração direta}

A demonstração direta é aquela que, por meio de uma sequência lógica, parte
das hipóteses e com argumentações coerentes concluímos a veracidade da tese.

\begin{proposicao}{}{par}
  Se $x$ é par, então $x^2$ é par.
\end{proposicao}

\begin{proof}
  Essa proposição nos diz que se um número $x$ é par, então o seu quadrado
  também é par. Note que dizemos que um número inteiro é par, se ele tem o
  formato $2k$, com $k \in \Z$ e dizemos que um número inteiro é ímpar,
  se ele tem o formato $2k+1$, com $k \in \Z$. Vamos inicialmente
  destacar a hipótese e a tese dessa proposição:
  \[
    \text{Hipótese: } x \text{ é par}.
  \]
  \[
    \text{Tese: } x^2 \text{ é par}.
  \]
  Assim, partindo da tese, se $x$ é um número par, então $x=2k$, onde $k$ é algum
  número inteiro. Assim,
  \[
    x^2= xx= (2k)(2k)= 4k^2=2(2k^2)
  \]
  Chamando $p=2k^2$, e notando que $2k^2$ é um número inteiro, temos que
  $x^2=2p$, então podemos concluir que $x^2$ é, também, um número par.
\end{proof}

Agora vamos provar o famoso Teorema de Pitágoras

\begin{teorema}{Teorema de Pitágora}{teo:teorema-pitagoras-novamente}
  Se $T$ é um triângulo retângulo
  de catetos $b$ e $c$ e hipotenusa $a$, então $a^2=b^2+c^2$.
\end{teorema}

\begin{figure}[h]
  \centering
  \includegraphics{teo-de-pitagoras}
  \caption{Triângulo $T$}
  \label{pita}
\end{figure}

\begin{figure}[h]
  \centering
  \includegraphics{teorema-de-pitagoras-quadrado}
  \caption{Quadrado formado pela junção de 4 triângulos retângulos iguais.}
  \label{pita1}
\end{figure}

\begin{proof}
  Aqui já sabemos que as hipótese são que $T$ é um triângulo retângulo de
  catetos $b$ e $c$ e hipotenusa $a$ e a tese, que é o resultado que queremos
  chegar é  $a^2=b^2+c^2$. Logo, considere um triângulo retângulo T conforme a
  Figura~\ref{pita}:

  A partir disto, vamos pegar 4 triângulos iguais os da Figura~\ref{pita} e
  posicionar estes triângulos de modo que formem um quadrado de lado $b+c$.
  Note que formamos também um quadrado interno de lados cujos lados são as
  hipotenusas $a$ dos triângulos, conforme ilustra a Figura~\ref{pita1}.

  Podemos calcular a área desse quadrado formado de duas formas. A primeira é
  notando que o lado desse quadrado é $b+c$, assim:
  \begin{equation} \label{eqpi}
    A_{\text{quadrado}} = (b+c)^2=b^2+2bc+c^2
  \end{equation}
  Por outro lado, podemos calcular a área desse quadrado maior somando a
  área dos 5 objetos que o compõem: 4 triângulos retângulos iguais e
  um quadrado menor de lado $a$.
  Ora,
  \[
    A_{\text{triangulo}} = \frac{bc}{2}
  \]
  \[
    A_{\text{quadrado menor}} = a^2
  \]
  Assim,
  \begin{equation}\label{eqpi2}
    A_{\text{quadrado}}
    = 4\frac{bc}{2}+a^2
    = 2bc+a^2
  \end{equation}
  Igualando as expressões obtidas em (\ref{eqpi}) e (\ref{eqpi2}) chegamos que:
  \[
    b^2+2bc+c^2= 2bc+a^2
  \]
  organizando os termos e fazendo os cancelamentos necessários, obtemos:
  \[
    a^2=b^2+c^2.
  \]
\end{proof}

%------------------------------------------------------------------------------%
\subsection{Demonstração por contraposição}

Sabemos que as proposições $P \Rightarrow Q$ e $\sim Q \Rightarrow \sim P$ são
equivalentes, ou seja,
$P \Rightarrow Q \Leftrightarrow \sim Q \Rightarrow \sim P$.
Nesse caso então iremos demonstrar de forma direta a sentença
$\sim Q \Rightarrow \sim P$,
que implica na demonstração de $P \Rightarrow Q$. Vamos
demonstrar as proposições que seguem:

\begin{proposicao}{}{}
  Se $X^2$ é par, então $X$ é par.
\end{proposicao}

\begin{proof}
  Inicialmente é importante observar, que nesse caso fica difícil fazer a
  demonstração direta (Verifique!). Vamos demonstrar a veracidade dessa
  afirmação usando a contrapositiva. Nesse caso, temos:
  \[
    P\colon X^2 \text{ é par}.
  \]
  \[
    Q\colon X \text{ é par}.
  \]
  Sendo assim,
  \[
    \sim Q\colon X \text{ é ímpar}.
  \]
  \[
    \sim P\colon X^2 \text{é ímpar}.
  \]
  Assumindo que $X$ é ímpar, temos que $X=2k+1$, $k \in \Z$. Assim,
  \[
    X^2
    = (2k+1)^2
    = 4k^2+4k+1
    = 2(2k^2+2k)+1
  \]
  Chamando $p=2k^2+2k$ e notando que $2k^2+2k$ é um número inteiro, temos que
  $X^2=2p+1$, que é o mesmo que dizer que $X^2$ é ímpar.
\end{proof}

%------------------------------------------------------------------------------%
\subsection{Demonstração por redução ao absurdo ou por contradição}

O método de redução ao absurdo, utiliza a negação da tese para chegar a uma
contradição da hipótese ou de alguma verdade matemática e é baseado na
\emph{lei do terceiro excluído} que diz que uma afirmação que não pode ser
falsa, deverá ser consequentemente verdadeira. Vamos ver na prática como essa
técnica funciona na prática provando a demonstração que segue:

\begin{proposicao}{}{}
  Se $a$ e $b$ são números reais não-negativos, então
  \(\frac{a+b}{2} \geq   \sqrt{ab}\).
\end{proposicao}

\begin{proof}
  Inicialmente, vamos destacar a hipótese e a tese:
  \[
    \text{Hipótese: } a, b \in \R_{+}
  \]
  \[
    \text{Tese: } \frac{a+b}{2} \geq \sqrt{ab}
  \]
  Começamos a demonstração negando a tese até chegar em algum absurdo.
  Assim, vamos supor que \(\frac{a+b}{2} < \sqrt{ab}\).
  Desta maneira
  \begin{align*}
    \frac{a+b}{2} & < \sqrt{ab} \\
    a+b & < 2 \sqrt{ab} \\
    a-2\sqrt{ab}+b & <0 \\
    (\sqrt{a}-\sqrt{b})^2 & <0 \\
  \end{align*}
  Deste modo essa conclusão é um absurdo, já que nenhum número real ao quadrado
  pode ser negativo. Portanto,
  \(\frac{a+b}{2} < \sqrt{ab}\) é falsa e
  \(\frac{a+b}{2} \geq \sqrt{ab}\) é verdadeira.
\end{proof}

Vamos usar a redução ao absurdo para provar a seguinte proposição:

\begin{proposicao}{}{}
  $\sqrt{2}$ é irracional.
\end{proposicao}

\begin{proof}
  Vamos negar a tese, ou seja, vamos supor que $\sqrt{2}$ é racional. Isso
  significa dizer que podemos escrever $\sqrt{2}$ na forma de fração
  irredutível, isto é, $\sqrt{2}=\sfrac{p}{q}$, com $p$ e $q$ primos entre si
  (não possui divisores comuns além do 1). Dessa maneira
  \[
    \sqrt{2}=\frac{p}{q}
    \quad\Rightarrow\quad 2=\frac{p^2}{q^2}
    \quad\Rightarrow\quad p^2=2q^2
  \]
  onde a última igualdade nos diz que $p$ é um número par, em outras palavras
  $p=2k$, $k \in \Z$. Assim, retornando para a última igualdade
  \[
    p^2=2q^2
    \quad\Rightarrow\quad (2k)^2=2q^2
    \quad\Rightarrow\quad 4k^2=2q^2
    \quad\Rightarrow\quad q^2=2k^2
  \]
  Dessa maneira concluímos que $q$ também é um número par. Ora, sendo $p$ e $q$
  números pares, então $2$ é um divisor comum de $p$ e $q$, o que é um absurdo,
  visto que estamos supondo de $p$ e $q$ são primos entre si. Portanto, é um
  absurdo supor que $\sqrt{2}$ é racional e sendo assim, $\sqrt{2}$ é
  irracional.
\end{proof}

%------------------------------------------------------------------------------%

%------------------------------------------------------------------------------%
\section{Exercícios}

\begin{exercises*}

%------------------------------------------------------------------------------%
\begin{exercise}
  Nos Teoremas e proposições a seguir, destaque a hipótese e a tese.
  \begin{alphalist}
    \item Se $a,b$ e $c$ são lados de um triângulo e $\hat{A}$ o ângulo
      oposto ao lado $a$, então $a^2 = b^2 + c^2 - 2bc \cos \hat{A}$
    \item Toda função diferenciável em $x=a$ é contínua em $x=a$. (Ou de
      modo equivalente: Se $f$ é uma função diferenciável então $f$ é uma
      função contínua em $x=a$.)
    \item Se $T$ é um triângulo equilátero então $Q$ é um triângulo
      isósceles (Ou de modo equivalente : Todo triângulo equilátero é
      isósceles).
    \item Se $x=1$, então $x^2=1$
  \end{alphalist}

  \begin{answer}
    \begin{alphalist}
      \item Hipótese: $a$, $b$ e $c$ são lados de um triângulo e $\hat{A}$
            é o ângulo oposto ao lado $a$ \\
            Tese: $a^2 = b^2 + c^2 - 2bc \cos \hat{A}$
      \item Hipótese: $f$ é uma função diferenciável em $x=a$ \\
            Tese: $f$ é uma função contínua em $x=a$
      \item Hipótese: $T$ é um triângulo equilátero \\
            Tese: $T$ é um triângulo isósceles
      \item Hipótese: $x=1$ \\
            Tese: $x^2=1$
    \end{alphalist}
  \end{answer}
\end{exercise}

%------------------------------------------------------------------------------%
\begin{exercise}
  Para as sentenças a seguir, dê a recíproca e a contrapositiva das
  proposições que seguem:
  \begin{alphalist}
    \item Se $k$ é um número par primo então $k=2$.
    \item Toda função diferenciável em $x=a$ é contínua em $x=a$
    \item Se $T$ é um triângulo equilátero então $T$ é um triângulo
      isósceles (Ou de modo equivalente : Todo triângulo equilátero é
      isósceles).
    \item Se $x=1$, então $x^2=1$
  \end{alphalist}

  \begin{answer}
    \begin{alphalist}
      \item Recíproca: Se $k=2$, então $k$ é um número par primo \\
        Contrapositiva: Se $k \neq 2$, então $k$ não é um número par primo
      \item Recíproca: Se $f$ é contínua em $x=a$,
        então $f$ é diferenciável em $x=a$ \\
        Contrapositiva: Se $f$ não é contínua em $x=a$,
        então $f$ não é diferenciável em $x=a$
      \item Recíproca: Se $T$ é isósceles, então $T$ é equilátero \\
        Contrapositiva: Se $T$ não é isósceles, então $T$ não é equilátero
      \item Recíproca: Se $x^2=1$, então $x=1$ \\
        Contrapositiva: Se $x^2 \neq 1$, então $x \neq 1$
    \end{alphalist}
  \end{answer}
\end{exercise}

%------------------------------------------------------------------------------%
\begin{exercise}
  A recíproca da proposição apresentada na questão anterior, item $e$ é
  verdadeira? Justifique.

  \begin{answer}
    Não. A recíproca afirma que se $x^2=1$ então $x=1$.
    Contudo, se $x = -1$, temos que $x^2 = 1$.
    Assim,  $x=\pm 1 \neq 1$, portanto a recíproca é falsa
  \end{answer}
\end{exercise}


%------------------------------------------------------------------------------%
\begin{exercise}
  Demonstre, de duas maneiras diferentes, as seguintes proposições. Em
  cada item, você deve escolher dois tipos de demonstração para fazer
  dentre: demonstração direta, redução ao absurdo e usando a
  contrapositiva.
  \begin{alphalist}
    \item A soma de dois números pares é sempre um número par.
    \item A soma de dois números ímpares é sempre um número par.
    \item O produto de dois números ímpares é sempre um número ímpar.
    \item Se $x^2$ é ímpar, então $x$ é ímpar.
    \item Se $x+y$ é divisível por $2$ então $x-y$ é divisível por $2$.
  \end{alphalist}

  \begin{answer}
    \begin{alphalist}
        \item
            Demonstração Direta:
            Sejam $n = 2k$ e $m = 2j$.
            A soma é $n + m = 2k + 2j = 2(k + j)$, que é par \\
            Redução ao Absurdo: Suponha que $2k + 2j = 2p + 1$.
            Então $2(k + j - p) = 1$, o que é uma contradição (par igual a ímpar)

        \item
            Demonstração Direta:
            Sejam $n = 2k + 1$ e $m = 2j + 1$.
            A soma é $(2k + 1) + (2j + 1) = 2(k + j + 1)$, que é par \\
            Redução ao Absurdo:
            Suponha que $2k+1 + 2j+1 = 2p + 1$.
            Então $2(k + j - p +1) = 1$, o que é uma contradição (par igual a ímpar)

        \item
            Demonstração Direta:
            Sejam $n = 2k + 1$ e $m = 2j + 1$.
            O produto é $4kj + 2k + 2j + 1 = 2(2kj + k + j) + 1$, que é ímpar \\
            Redução ao Absurdo:
            Suponha que o produto seja par ($2p$).
            Então $2(2kj+k+j) + 1 = 2p \Rightarrow 1 = 2(p - 2kj - k - j)$,
            o que é uma contradição (par igual a ímpar)

        \item
            Usando a Contrapositiva:
            Se $x=2k$, então $x^2 = 4k^2 = 2(2k^2)$, que é par.
            Provada a contrapositiva, a original é verdadeira \\
            Redução ao Absurdo:
            Suponha que $x$ é par.
            Então $x^2 = (2k)^2 = 2(2k^2)$ que é par,
            o que contradiz a hipótese de $x^2$ ser ímpar

        \item
            Demonstração Direta:
            Se $x + y = 2k$, então $x - y = (x + y) - 2y = 2k - 2y = 2(k - y)$,
            que é divisível por 2 \\
            Redução ao Absurdo:
            Suponha que $x - y$ seja ímpar ($2j+1$).
            Somando com $x+y=2k$, teríamos $2x = 2(k+j)+1$,
            o que é uma contradição (par igual a ímpar)
    \end{alphalist}
  \end{answer}
\end{exercise}


%------------------------------------------------------------------------------%
\begin{exercise}
  Demonstre, usando redução ao absurdo, que $\sqrt{3}$ é um número
  irracional.

  \begin{answer}
    Redução ao absurdo: Suponha $\sqrt{3} = \frac{p}{q}$ irredutível.
    Então $3 = \frac{p^2}{q^2} \Rightarrow p^2 = 3q^2$.
    Logo, $p = 3k$ \\
    Substituindo: $(3k)^2 = 3q^2 \Rightarrow 3k^2 = q^2$.
    Logo, $q$ também é múltiplo de 3, contradizendo a irredutibilidade
  \end{answer}
\end{exercise}


%------------------------------------------------------------------------------%
\begin{exercise}
  Todas as proposições abaixo são falsas. Dê um contra-exemplo que
  invalide cada uma das afirmações.
  \begin{alphalist}
    \item Em todo triângulo, de lados $a,b$ e $c$, e com $a$ sendo o maior
      lado, vale que $a^2=b^2+c^2$.
    \item Se $a$ e $b$ são números não-negativos então $\sqrt{a+b}=
      \sqrt{a}+ \sqrt{b}$.
    \item Se $a$ é um número real qualquer, então $a^2>a$.
    \item Todo número natural ímpar é primo.
    \item Todo triângulo é retângulo.
  \end{alphalist}

  \begin{answer}
    \begin{alphalist}
      \item Contraexemplo: Triângulo de lados $a=5, b=4, c=2$.
            Temos $5^2 \neq 4^2 + 2^2$
      \item Contraexemplo: $a=9, b=16$.
            Temos $\sqrt{9+16}=5$, mas $\sqrt{9}+\sqrt{16}=7$
      \item Contraexemplo: $a=0.5$.
            Temos $0.5^2 = 0.25$, e $0.25 < 0.5$
      \item Contraexemplo: $n=9$ é ímpar,
            mas não é primo (é divisível por 3)
      \item Contraexemplo: Triângulo equilátero possui todos os
            ângulos iguais a $60^\circ$
    \end{alphalist}
  \end{answer}
\end{exercise}

\end{exercises*}

%------------------------------------------------------------------------------%


%------------------------------------------------------------------------------%
